\documentclass[12pt]{scrartcl}
\usepackage[a4paper,margin=3.0cm]{geometry}
\usepackage{../common/rlstudy}
\linespread{1.15}\selectfont

\title{Extended reading --- Week 03: Monte Carlo Methods}
\date{\today}
\begin{document}
\maketitle

\section*{How to use this list}
The first paper is the source of every stochastic-approximation
guarantee in RL. The second is the resolution of the first-visit vs
every-visit MC question. The third is the canonical reference for
per-decision IS (the variance-reduction trick that makes off-policy RL
tractable). The fourth is the modern reference for off-policy evaluation
with high-confidence bounds, which you should read \emph{before} deploying
any new policy. The fifth is the bridge from off-policy MC to deep learning.

\section{Robbins \& Monro (1951) --- A stochastic approximation method}
\textit{Ann.\ Math.\ Stat.\ 22, 400--407.}

\paragraph{Summary.} Introduces the stochastic-approximation iteration
$\theta_{n+1} = \theta_n + \alpha_n (Y_n - f(\theta_n))$ for finding the root
of an unknown regression function. Proves convergence in mean-square if
$\sum \alpha_n = \infty$ and $\sum \alpha_n^2 < \infty$ --- the celebrated
``RM conditions''.

\paragraph{Why it matters for this week.} The incremental sample-average
$Q_{n+1} = Q_n + \tfrac{1}{n}(R_n - Q_n)$ is exactly the Robbins--Monro
algorithm with $\alpha_n = 1/n$. The convergence proofs of every
``learning rate goes to zero'' RL algorithm (SARSA, Q-learning, TD($\lambda$)
on-policy) hinge on Robbins--Monro.

\begin{quote}\itshape
``Convergence in mean-square is established under conditions that include
$\sum \alpha_n = \infty$, ensuring the iteration moves forever, and
$\sum \alpha_n^2 < \infty$, ensuring the noise averages out.''
\end{quote}
``Move forever, but slow down quickly enough that noise averages out'' is the
single most important piece of intuition for learning-rate selection.

\section{Singh \& Sutton (1996) --- Reinforcement learning with replacing
eligibility traces}
\textit{Machine Learning 22, 123--158.}

\paragraph{Summary.} Compares first-visit and every-visit MC and shows the
bias of every-visit is $O(1/N)$. Introduces \emph{replacing} eligibility
traces (overwrite, not accumulate) and proves they correspond to first-visit
MC in the $\lambda=1$ limit. Sets up the unification of MC and TD via
$\lambda$ that Week 4 builds on.

\paragraph{Why it matters for this week.} It is the reason ``first-visit''
shows up as the default in MC reference implementations. Read \S 3 of the
paper for the cleanest treatment of the bias issue. The same paper also
explains the otherwise mysterious choice ``replacing'' vs ``accumulating''
traces in classical TD($\lambda$).

\begin{quote}\itshape
``First-visit MC produces independent samples of the return per state per
episode; every-visit MC produces dependent samples but does not waste data.
The bias of every-visit vanishes as $1/N$ in the number of episodes.''
\end{quote}
The ``does not waste data'' is the practical case for every-visit; the
``independent samples'' is the theoretical case for first-visit. In most
deep-RL implementations the distinction is academic --- the variance is
dominated by other factors.

\section{Precup, Sutton \& Singh (2000) --- Eligibility traces for off-policy
policy evaluation}
\textit{ICML 2000.}

\paragraph{Summary.} Defines per-decision importance sampling for MC and TD
methods. Proves that the per-decision estimator has the same expectation as
the trajectory-level OIS but lower variance whenever the immediate rewards
depend on the action at time $t$ but not on later actions. Generalises to
TD($\lambda$) with off-policy targets.

\paragraph{Why it matters for this week.} If you remember one off-policy
trick, remember per-decision IS. The trajectory-level estimator is unbiased
but exponentially noisy in the horizon $T$; per-decision IS keeps the
unbiasedness but only multiplies by the importance ratio \emph{up to} time
$t$ for each reward $R_{t+1}$. The variance reduction is dramatic in any
long-horizon setting.

\begin{quote}\itshape
``Per-decision importance sampling reduces variance without introducing bias
by noticing that the reward at time $t$ is independent of the actions taken
after time $t$ given the policy.''
\end{quote}
That single observation underpins every modern off-policy estimator. It is
also what makes truncated trust-region methods possible (PPO, GRPO).

\section{Thomas \& Brunskill (2016) --- Data-efficient off-policy policy
evaluation for reinforcement learning}
\textit{ICML 2016.}

\paragraph{Summary.} Introduces the \emph{doubly-robust} estimator for RL:
combine importance sampling with a learned reward / value model so the
estimate is consistent if \emph{either} the model or the propensity is
correct. Provides finite-sample concentration bounds and high-confidence
off-policy evaluation: a $(1-\delta)$ lower bound on $V^\Policy$ that can be
inspected before deployment.

\paragraph{Why it matters for this week.} Read this before deploying any
candidate policy to production. The DR estimator is the right default in any
applied setting with both logs and a reasonable reward model. The
high-confidence lower bound (Section 4.3) is the deployment gate: only ship
a policy whose lower bound exceeds the production policy's upper bound.

\begin{quote}\itshape
``A safe policy improvement scheme: never deploy a policy whose 95\% lower
bound on value is below the current policy's empirical value. Doubly-robust
estimation provides the bound that makes this protocol tight.''
\end{quote}
``Never deploy unless the lower bound dominates'' is the cleanest production
rule in RL. It is also the rule that, if synthetic insurance / ATS had it on day one,
would have prevented several historical regressions.

\section{Joachims, Swaminathan \& de Rijke (2018) --- Deep learning with
logged bandit feedback}
\textit{ICLR 2018.}

\paragraph{Summary.} Trains a deep neural-net policy on logged bandit
feedback using the SNIPS estimator as a counterfactual risk objective.
Adds \emph{propensity overfitting} controls (self-normalisation, weight
clipping) to keep deep counterfactual training stable. Bridges classical
off-policy MC to modern offline-RL methods like CQL.

\paragraph{Why it matters for this week.} It is the deep-learning analogue of
your week-3 workbook exercise: estimate the value of a target policy from
logs and \emph{optimise} the policy against that estimate. The propensity-
overfitting analysis is exactly the safety concern of the lecture.

\begin{quote}\itshape
``Self-normalisation defeats propensity overfitting --- the failure mode where
the policy learns to put mass on rare-but-high-weight logged examples. We
demonstrate this failure empirically and show that SNIPS-based training
avoids it.''
\end{quote}
``Propensity overfitting'' is the deep-learning name for the OIS variance
explosion. Knowing both names lets you talk to both communities.

\section*{Reading exercises}
\begin{enumerate}\setlength\itemsep{4pt}
  \item \textbf{Robbins--Monro, optimal $\alpha_n$.} Robbins--Monro prove
        convergence under any $\alpha_n$ satisfying the two summability
        conditions. The optimal rate $\alpha_n = c/n$ achieves $1/\sqrt{n}$
        mean-square error. Read Polyak \& Juditsky (1992, \emph{SIAM J.\
        Control}) on iterate averaging --- show that averaging the iterates
        of a slower-than-optimal step size recovers $1/\sqrt{n}$ asymptotic
        variance.
  \item \textbf{Singh--Sutton, finite-sample comparison.} Run first-visit and
        every-visit MC on a chain MDP with $N\in\{10, 100, 1000\}$ episodes.
        Plot the bias and variance of each as a function of $N$. Verify the
        $1/N$ bias claim for every-visit.
  \item \textbf{Precup--Sutton--Singh, per-decision IS on chains.} Implement
        per-decision IS on a 5-state chain. Compare its variance with
        trajectory-level OIS at $T=20$. Verify that the variance ratio scales
        like $T \cdot \E[\rho^2]$ vs $\E[\rho_\tau^2]$.
  \item \textbf{Thomas--Brunskill, doubly-robust on Mountain Car.} Train a
        small reward model on logs of a sub-optimal policy and use it in the
        DR estimator. Compare the high-confidence lower bound of DR with
        WIS on a Mountain-Car-like environment.
  \item \textbf{Joachims et al., BanditNet on synthetic clicks.} Reproduce
        Algorithm 1 from the paper on the synthetic click dataset in their
        Section 5. Compare with a naive cross-entropy classifier; how much
        does propensity overfitting cost the naive approach?
\end{enumerate}

\section*{Connections to later weeks}
\begin{itemize}\setlength\itemsep{2pt}
  \item \emph{Week 4 (TD)} replaces the MC return by a bootstrap target;
        IS still applies but per-decision.
  \item \emph{Week 11 (contextual bandits + OPE)} unifies SNIPS, DR, SWITCH,
        and MIPS in a single contextual-bandit benchmark.
  \item \emph{Week 14 (offline RL)} extends doubly-robust to full RL
        (Jiang \& Li 2016) and is the foundation for conservative offline
        Q-learning.
  \item \emph{Week 19 (DPO and friends)} uses an importance-weighted
        contrastive loss directly on preferences --- a direct descendant
        of the off-policy MC ideas of this week.
\end{itemize}
Monte Carlo is conceptually simple but its lessons (coverage, importance
weights, doubly-robust safety) recur in every off-policy method that follows.

\section*{Sourced from}
\verb|research/bandits_matching_sota.md| \S 3 and Sutton \& Barto Chapter 5.
For a deeper dive into OPE see Voloshin, Le, Jiang, Yue \emph{Empirical
Study of Off-Policy Policy Evaluation for RL} (NeurIPS 2021 benchmarks).

\end{document}
